\begin{theorem}[MetricSpace L10 sibling lattice route]
\label{thm:metricspace-l10-sibling-lattice-route}
Every $\MetricSpaceUp$ real-distance consumer used by the current L10 Real
and Completion route factors through the finite sibling surfaces already
present in the metric child files:
\autoref{ch:concrete-instances-locatedreal-namecert},
\autoref{ch:concrete-instances-s1-namecert},
\autoref{ch:concrete-instances-real-namecert},
\autoref{ch:concrete-instances-streamname-namecert},
\autoref{ch:concrete-instances-regseqrat-namecert}, and
\autoref{ch:concrete-instances-dyadicratcore-namecert}. The lattice route is
the displayed finite read through LocatedReal endpoints, SOne component
routes, the Real public-distance certificate, StreamName finite windows,
RegSeqRat regular-source rows, and DyadicRatCore endpoint and budget ledgers.

The route introduces no host metric function, quotient real equality,
selected limit, open-cover compactness, or metric-completion datum.
\end{theorem}
\leanchecked{BEDC.Derived.MetricUp.MetricspaceL10SiblingLatticeRoute\_finite\_packet}
\begin{proof}
First place the located-real side on its finite source boundary. Apply
\autoref{thm:metricspace-locatedreal-distance-source-boundary}. The distance
consumer is then read only through accepted $\LocatedRealUp$ packets, their
$\StreamNameUp$ schedules, $\RegSeqRatUp$ source windows, and
$\DyadicRatCore$ endpoint and radius ledgers, together with visible
$\BHist$, $\hsame$, $\Cont$, $\Pkg$, and $\NameCert$ rows.

Next place the circle-side metric consumer on its finite component route.
Apply \autoref{thm:metricspace-sone-distance-source-boundary}. This keeps the
$\SOneUp$ read on the two $\RealUp$ component rows, the unit-equation row,
componentwise classifiers, finite-window stream and regular-source rows, and
the same visible provenance surface.

Now attach the Real-facing certificate by
\autoref{thm:metricspace-real-distance-scoped-certificate}. Its certificate
is exactly the real-distance carrier, the public metric laws, the
RealAlgOrder positive-distance correspondence, and the apartness triangle
boundary. Therefore the metric consumer reaches the L10 Real face through the
same finite packet rather than through a separate real equality or completion
coordinate.

Finally apply \autoref{thm:metricspace-public-law-package}. The exported
metric laws are symmetry, identity boundary, triangle-budget composition, and
carrier transport over the displayed packet. Since the LocatedReal and SOne
source boundaries already factor through the named Real, StreamName,
RegSeqRat, and DyadicRatCore sibling faces, the assembled L10 route is the
finite sibling lattice listed in the statement. None of the cited rows
contains a host metric function, quotient real equality, selected limit,
open-cover compactness, or metric-completion datum, so none is introduced by
the factorization.
\end{proof}
\leanchecked{BEDC.Derived.MetricUp.MetricSpaceL10SiblingLatticeRoute}
